\scanpage{5}
\mspage{2}

\deleted{d) Soit}

\[
i:Y\lhook\joinrel\longrightarrow S
\]

est une immersion fermée, de complémentaire $U$

\diagram{\begin{tikzcd}[ampersand replacement=\&]
Y\arrow[r,hook,"i"] \& S \& U\arrow[l,hook',"j"']
\end{tikzcd}}

de sorte que $j$ est une imm. ouverte, alors $A(Y)$ s'identifie \inserted{par le foncteur $i_{*}$} à la sous-catégorie pleine de $A(S)$ formée des $F\in\mathrm{Ob}\,A(S)$ tels que $F|U=e_{U}$ (préfaisceau final). \unsure{De $=$|De même}, si l'objet initial $\unsure{\varphi|\emptyset}$ existe dans $A$, et les $\underset{-}{\lim}$ pertinentes, alors $A(U)$ s'identifie par $j_{!}$ à la sous-catégorie pleine de $A(S)$ formée des $F$ tels que $F|Y=\varphi_{Y}$ (préfaisceau initial, \deleted{\unsure{j}} prenant partout la valeur $\varphi$.

Si \deleted{il existe} $A$ est $=$ \uline{pointée} ($e=\varphi$), alors on a pour tout $F\deleted{\unsure{\in\mathrm{Ob}}}\,A(S)$ un \deleted{\unsure{d}}carré à la fois cartésien et cocartésien, fonctoriel en $F$

\diagram{\begin{tikzcd}[ampersand replacement=\&]
\& \& F\arrow[dr,"\text{adj}"] \& \\ (1) \& j_{\deleted{\unsure{*}}!}j^{*}F\arrow[ur,"{\text{(adj.)}}"]\arrow[dr] \& \substack{\text{cart}\\\text{et}\\\text{cocart}} \& i_{*}i^{*}F \\ \& \& e\arrow[ur] \&
\end{tikzcd}}

Dans le cas $A$ abélienne, ce diagramme prend la forme d'une suite exacte

\[
0\longrightarrow j_{!}j^{*}F\unsure{-|\longrightarrow}F-i_{*}i^{*}F\longrightarrow0
\]

e) Mêmes notations pour $i$ et $j$. On suppose que $A$ a les $\varprojlim$ pertinentes, donc objet final. On a alors un carré cartésien fonctoriel en $F\in\mathrm{Ob}\,A(S)$

\diagram{\begin{tikzcd}[ampersand replacement=\&]
\& F\arrow[d,"\text{adj}"']\arrow[dr,phantom,"\text{cart}" description] \& i_{*}i^{!}F\arrow[l]\arrow[d] \\ (\unsure{1|2}) \& \deleted{\unsure{i_{*}i}}j_{*}j^{*}F \& e\arrow[l]
\end{tikzcd}}

\unsure{i.e.|i.e.,} $i^{!}F$ s'identifie par le foncteur pl. fid $i_{*}$ au \uline{noyau} de $F\longrightarrow j_{*}j^{*}F$. Dans le cas abélien cel\unsure{a} prend la forme d'une suite exacte

\[
0\longrightarrow i_{*}i^{!}F\longrightarrow F\longrightarrow j_{*}j^{\unsure{*|\prime}}F\ .
\]

% archive sheet number: 4
